\title{Parameter-Efficient Orthogonal Finetuning via Factorization}

\begin{abstract}
The widespread adoption of large foundation models underscores their utility, yet the exorbitant costs associated with training them anew are a significant barrier. Consequently, devising effective strategies for adapting these models to specific downstream applications has become increasingly crucial. This work investigates a theoretically grounded finetuning framework, Orthogonal Finetuning (OFT), as a means of downstream adaptation. While OFT exhibits robust generalization capabilities, it remains parameter-intensive because orthogonal matrices are inherently high-dimensional and thus costly to learn. To overcome this, we analyze OFT through the lens of information transmission and articulate several important criteria for achieving greater parameter efficiency. Drawing inspiration from the efficiency of the Cooley-Tukey fast Fourier transform in facilitating information transfer, we propose a streamlined orthogonal parameterization based on butterfly structures. We leverage this parameterization within OFT, leading to a new, highly parameter-efficient finetuning strategy: Orthogonal Butterfly (BOFT). By framing OFT as a particular instance within this broader paradigm, BOFT introduces a unified and flexible orthogonal finetuning framework. We empirically validate BOFT by extensively adapting state-of-the-art vision transformers, large language models, and text-to-image diffusion architectures to a variety of vision and language tasks.
\end{abstract}

\section{Introduction}

The latest foundation models, such as ChatGPT~\cite{brown2020language,chen2021evaluating} and Stable Diffusion~\cite{rombach2022high}, exemplify the superior generalization properties achievable by scaling up model sizes. However, this progress has been accompanied by an explosive growth in model parameter counts—GPT-3, for example, comprises approximately 175 billion parameters—which presents substantial barriers for training models from the ground up. As a result, advancing the field now critically depends on our ability to repurpose existing foundation models without incurring the cost of complete retraining. In practice, this necessitates efficient methodologies for transferring models to new tasks. Currently, three predominant adaptation strategies exist: (1) \emph{model finetuning}~\cite{hulora2022,zaken2022bitfit,guo2020parameter,raffel2020exploring,qiu2023controlling,zhang2022adaptive,chavan2023one}, which preserves the underlying architecture while selectively finetuning subsets of parameters; (2) \emph{adapter tuning}~\cite{rebuffi2017learning,houlsby2019parameter,pfeiffer2020adapterfusion,lin2020exploring,he2021towards}, which integrates additional trainable modules into the original model; and (3) \emph{prompt tuning}~\cite{li2021prefix,lester2021power}, which prepends trainable token prefixes to the model's input. Of these, model finetuning stands out for its straightforward implementation, notable effectiveness, and most importantly, for introducing no additional latency during inference.

Model finetuning is fundamentally guided by the idea that the finetuned model should remain close to its pretrained counterpart according to specific metrics, thereby retaining the knowledge obtained during pretraining. Conventional finetuning strategies typically employ a low learning rate, an empirical technique that implicitly constrains the deviation between the two models. Recognizing the organized structure inherent in weight matrices, a more systematic strategy aims to maintain the relational properties of these matrices, particularly preserving the pairwise angular relationships among neurons. This motivation underpins the development of Orthogonal Finetuning~(OFT)~\cite{qiu2023controlling}, a finetuning paradigm in which all neurons within a single layer undergo transformation via a shared orthogonal matrix, thereby guaranteeing that the angles between neuron vectors are maintained throughout finetuning. While OFT has shown notable effectiveness and robust convergence when applied to finetuning text-to-image diffusion models~\cite{qiu2023controlling}, its practicality is hindered by the substantial parameter count necessitated by high-dimensional orthogonal matrices. To alleviate this, OFT incorporates a block-diagonal parameterization to achieve greater parameter efficiency. Nonetheless, this comes with certain limitations—namely, the enforced block-wise sparsity of the orthogonal matrix restricts flexibility, as each block operates independently and follows a fixed sparsity pattern. As a result, this structure, while reducing the number of parameters, may introduce unintended inductive biases (such as limiting the expressivity of the orthogonal transformation and impeding the approximation of standard linear mappings). Furthermore, there remains an open question as to how one might determine an optimal block-wise sparsity pattern.

The central challenge in this context is constructing a dense orthogonal matrix while maintaining a compact parameterization. Ordinarily, a dense orthogonal matrix of dimension $d$ would necessitate $\mathcal{O}(d^2)$ parameters, posing a significant obstacle. In response, we introduce an alternative methodology, whereby a dense orthogonal matrix is composed via a product of multiple factorized sparse matrices. This strategy stems from the observation that parameter count can be curtailed by exchanging memory requirements for additional computational cost. Since our orthogonal matrix is expressed as a multiplication of sparse factors, the necessary multiplications must be performed repeatedly throughout training. To elucidate this matrix factorization, we recast the problem of constructing a dense orthogonal matrix as one of information transmission. In this framing, synthesizing a dense orthogonal matrix via sparse multiplicative factors becomes analogous to propagating information over a grid-like graph structure. This interpretation facilitates the exploration of various sparse matrix factorization schemes, each limiting the number of adjustable parameters yet still capable of representing dense matrices effectively.

To enhance parameter efficiency within this information transmission paradigm, we look to the butterfly structure central to the Cooley-Tukey fast Fourier transform algorithm~\cite{cooley1965algorithm}, renowned for its capacity to allow efficient communication between disparate nodes~\cite{klasing1994broadcasting}. The butterfly graph intrinsic to the Cooley-Tukey procedure provides a principled mechanism for sparse matrix factorization, which we refer to as butterfly factorization. Through this construction, a dense $d$-dimensional matrix can be represented as a product of $\mathcal{O}(\log d)$ sparse matrices, with each matrix possessing only $\mathcal{O}(d)$ non-zero elements. By constraining each sparse factor to be orthogonal, we obtain a highly efficient orthogonal matrix representation that requires merely $\mathcal{O}(d\log d)$ parameters—an immense reduction relative to the conventional approach. Building upon this streamlined orthogonal parameterization, we introduce a new finetuning technique, Orthogonal Butterfly (BOFT), that emphasizes parameter efficiency. OFT is encapsulated within BOFT as a specific case, making BOFT a general orthogonal finetuning schema. Both the butterfly and block-diagonal designs possess a shared sparsity characteristic, injecting structured sparsity into the orthogonal matrices and thus decreasing the total number of parameters that must be learned. In juxtaposition, one may compare our approach to LoRA~\cite{hulora2022}, which leverages low-rank matrix structures for efficient parameterization; notably, both low-rank and sparse matrices are prominent classes of structured matrices~\cite{candes2011robust} exploited to optimize parameter usage.

In contrast to the block-diagonal configuration leveraged by OFT to balance the trade-off between expressivity and regularity, BOFT utilizes the butterfly structure to achieve a more continuous interpolation between matrices belonging to the full orthogonal group (maximizing expressivity) and identity matrices (promoting regularity). This formulation allows for the isolation of a more compact hypothesis space within the orthogonal group, which is favorable for downstream applications. Given that the butterfly structure underlies several efficient linear operators—such as the discrete Fourier transform and the discrete cosine transform—we posit that our structured, parameter-efficient methodology imparts an implicit inductive bias, enhancing model generalization while mitigating the risk of overfitting. This intuition is underpinned by the butterfly structure's capacity to easily represent various canonical linear transformations—a property not shared by the block-diagonal arrangement employed in OFT. Our principal contributions are as follows:

\begin{itemize}[leftmargin=*,nosep]
    \item We reframe the challenge of parameter efficiency in orthogonal finetuning via an innovative information transmission perspective. Specifically, we reinterpret the task of designing a parameter-efficient dense orthogonal matrix as an information routing problem mapped onto a grid-structured graph.
    \item Drawing inspiration from the butterfly architecture employed in the Cooley-Tukey algorithm, we introduce Orthogonal Butterfly—a parameter-efficient method for orthogonal finetuning—formulated within the context of our information transmission framework.
    \item We offer several theoretical justifications for BOFT’s ability to markedly reduce the count of learnable parameters, yet maintain robust expressivity and stable training dynamics. Owing to its matrix factorization approach, BOFT further enables an intriguing form of weight interpolation (refer to Figure~\ref{fig:boft_interpolation}).
    \item For the first time, we demonstrate the application of orthogonal finetuning across a variety of tasks that extend beyond the scope of controllable text-to-image generation~\cite{qiu2023controlling}, highlighting BOFT’s promise as a versatile model finetuning paradigm. Specifically, we evaluate BOFT on diverse adaptation challenges spanning computer vision and natural language processing, where it substantially surpasses current leading methods, affirming both its parameter efficiency and ability to generalize.
\end{itemize}

\section{Related Work}

\textbf{Parameter-efficient finetuning (PEFT)}. The rapid growth and increasing capabilities of foundation models (\emph{e.g.}, \cite{brown2020language,radford2021learning,rombach2022high,kirillov2023segment}) have intensified research interest in finetuning these models for specific downstream tasks with a focus on parameter efficiency~\cite{treviso2023efficient,ding2023parameter,lialin2023scaling}. Within the broad array of PEFT strategies~\cite{houlsby2019parameter,aghajanyan2020intrinsic,hulora2022,edalati2022krona,wang2022adamix,gheini2021cross,zaken2022bitfit,guo2020parameter,sung2021training,ansell2022composable,lester2021power,li2021prefix,vu2022spot,he2021towards,mao2021unipelt,karimi2021compacter,liu2022few,sung2022lst,chen2022parameter,jia2022visual,chen2022adaptformer,zhang2022neural,jie2023fact,lian2022scaling,luo2023towards}, reparameterization-based approaches~\cite{aghajanyan2020intrinsic,hulora2022,edalati2022krona} are most pertinent to our study. LoRA~\cite{hulora2022} introduces two low-rank matrices whose product is added to the pretrained weight matrix, delivering strong results on language-related tasks. While LoRA employs a fixed rank across all additional matrices, more recent approaches \cite{zhang2022adaptive,valipour2022dylora,zhang2023increlora} adaptively assign the rank at different network layers for better utilization of parameter resources. Because of its straightforwardness, low-rank weight reparameterization has gained widespread adoption~\cite{dettmers2023qlora,zi2023delta,chavan2023one}. Drawing motivation from the role of hyperspherical energy in representing generalization~\cite{liu2018learning,liu2021orthogonal}, \cite{qiu2023controlling} put forth orthogonal finetuning, a distinctive yet potent method for adapting text-to-image diffusion models. In this method, OFT learns an orthogonal transformation at the layer level, enabling stronger generalization and consistently more stable optimization than LoRA. However, OFT generally requires more trainable parameters than LoRA, indicating that further parameter efficiency would be valuable. Additionally, OFT's applicability beyond text-to-image diffusion models~\cite{qiu2023controlling} has not been extensively explored. By leveraging butterfly factorization, BOFT enhances the parameter efficiency of OFT, which, in turn, allows us to extend orthogonal finetuning to broader adaptation scenarios.

\textbf{Butterfly structures}. The radix-2 Cooley-Tukey algorithm transforms the $N$-point discrete Fourier transform recursively into two smaller $\frac{N}{2}$-point transforms, thereby generating a butterfly structure that can be expressed as the product of multiple sparse matrices (collectively termed a butterfly matrix). Butterfly matrices have previously served to parameterize orthogonal matrices to sidestep pivoting in Gaussian elimination and improve computational efficiency~\cite{parker1995random}, stabilize the learning of recurrent neural networks~\cite{jing2017tunable}, and facilitate kernel approximation~\cite{munkhoeva2018quadrature}. Fast linear transformation learning via butterfly parameterizations has been demonstrated in \cite{dao2019learning,dao2019kaleidoscope}. The application of butterfly matrices to efficiently train neural networks has been investigated in \cite{chen2022pixelated,dao2022monarch}. Moreover, butterfly structures play a role in efficient matrix-vector multiplication~\cite{michielssen1996multilevel,o2010algorithm}, approximating data-sparse matrices~\cite{li2015butterfly}, and enhancing network communication~\cite{klasing1994broadcasting,li2003linear,rajasingh2016transmission}. Unlike prior efforts, our work aims to exploit butterfly structures specifically to improve the parameter efficiency of OFT.

\section{An Information Transmission View on Orthogonal Finetuning}

We begin by outlining foundational aspects of OFT. In the context of finetuning pretrained weight matrices, OFT reformulates the new weight matrix as the multiplication of a trainable orthogonal matrix and the fixed pretrained matrix. Unlike LoRA, which incorporates an additive low-rank matrix for weight adaptation, OFT employs a multiplicative orthogonal matrix as the update mechanism. LoRA attains parameter efficiency via its low-rank construction, whereas the original OFT utilizes a block-diagonal orthogonal matrix structure~\cite{qiu2023controlling}; the parameter efficiency of OFT improves as the block sizes decrease. Figure~\ref{fig:comp_lora} visually contrasts these approaches. The rationale for using orthogonal transformations in OFT is to maintain the angular relationships between neurons~\cite{liu2018learning,liu2021orthogonal,qiu2023controlling}, thereby preserving the bulk of the pretrained model’s semantic information. Specifically, OFT involves optimizing an orthogonal matrix $\bm{R}\in\mathbb{R}^{d\times d}$ for a given pretrained linear layer $\bm{W}^0\in\mathbb{R}^{d\times n}$, changing the forward computation from $\bm{z}=(\bm{W}^0)^\top\bm{x}$ to $\bm{z}=(\bm{R}\bm{W}^0)^\top\bm{x}$, where $\bm{x}\in\mathbb{R}^{d}$ and $\bm{z}\in\mathbb{R}^{n}$ are the respective input and output vectors. For zero initialization, OFT sets $\bm{R}$ to be the identity matrix. To maintain orthogonality of $\bm{R}$ throughout finetuning, we adopt the Cayley parameterization as detailed in~\cite{qiu2023controlling,liu2021orthogonal}, specifically: $\bm{R}=(\bm{I}+\bm{Q})(\bm{I}-\bm{Q})^{-1}$, where $\bm{Q}$ is a skew-symmetric matrix, satisfying $\bm{Q}=-\bm{Q}^\top$. For improved parameter efficiency, the block-diagonal format is used, representing $\bm{R}$ as $\text{diag}(\bm{R}_1,\bm{R}_2,\cdots,\bm{R}_r)$, with each $\bm{R}_i\in\mathbb{R}^{b\times b}$ being a smaller orthogonal matrix such that $br=d$. However, this structural advantage imposes a limitation: it assumes that each neuron's dimensions (that is, columns of $\bm{W}^0$) can be partitioned into $r$ groups, each group being independently transformed by separate orthogonal matrices. Despite empirical support for this approach, the rationale behind such index-based grouping is tenuous, motivating interest in dense orthogonal matrices. Consequently, we pose the following: \emph{Is it possible to construct a dense orthogonal matrix while retaining parameter efficiency?}

In response to this inquiry, we suggest constructing a dense orthogonal matrix by multiplying several sparse orthogonal matrices. To offer a coherent and accessible framework for analyzing the sparsity structure in orthogonal matrix factorizations, we recast the generation of a dense orthogonal matrix within OFT as an information transmission paradigm. More precisely, forming a dense matrix $\bm{R}\in\mathbb{R}^{d\times d}$ as the product $\thickmuskip=2mu \medmuskip=2mu \bm{R}=\bm{B}_m\bm{B}_{m-1}\cdots\bm{B}_1$—where each $\bm{B}_i$ is a square matrix—can be interpreted as the process of passing information through a lattice of $d\times (m+1)$ nodes (see Figure~\ref{fig:info_trans}). This information-theoretic perspective is motivated by noting that a dense $d$-by-$d$ matrix encodes comprehensive connectivity between two groups of $d$ nodes. In $\bm{R}$, an entry $R_{ij}=0$ signifies that there is no transmission from the $j$-th input node to the $i$-th output node, whereas a nonzero $R_{ij}$ indicates possible information flow. Hence, expressing $\bm{R}$ as $\bm{B}_m\bm{B}_{m-1}\cdots\bm{B}_1$ can be interpreted as a sequence of information transfers governed by the sequence of graphs represented by each $\bm{B}_i$. The pathway of information thus starts with the structure imposed by $\bm{B}_1$ and progresses through to $\bm{B}_n$. As illustrated concretely in Figure~\ref{fig:info_trans}, consider the decomposition $\bm{R}=\bm{B}_5\bm{B}_4\bm{B}_3\bm{B}_2\bm{B}_1$, for which both the sparsity configurations and the resulting graph are depicted. The graph shown results from unraveling the process of matrix multiplication. Here, each $\bm{B}_i$ serves as a connectivity matrix linking nodes at the $i$-th layer to those at the $(i+1)$-th layer; more explicitly, the $(j_1,j_2)$ element of $\bm{B}_i$ represents whether there is a directed connection from the $j_2$-th node at the $i$-th level to the $j_1$-th node at level $(i+1)$ (a value of zero indicates absence of a connection). For the overall product $\bm{B}_5\bm{B}_4\bm{B}_3\bm{B}_2\bm{B}_1$ to result in a dense matrix, each starting node at the first layer must be capable of reaching any node at the sixth layer. By reducing to $\bm{R}=\bm{B}_4\bm{B}_3\bm{B}_2\bm{B}_1$—corresponding to paths from first-layer source nodes to fifth-layer target nodes—we observe that node $1$ cannot transmit information to node $3$. Thus, the pattern of zeros in $\bm{B}_4\bm{B}_3\bm{B}_2\bm{B}_1$ includes a zero at entry $(3,1)$. Similar relationships between the induced graphs and sparsity patterns hold for products such as $\bm{R}=\bm{B}_3\bm{B}_2\bm{B}_1$ and $\bm{R}=\bm{B}_2\bm{B}_1$. Broadly, in order for a matrix $\bm{R}\in\mathbb{R}^{d\times d}$ to be dense, the $m$ factor matrices $\bm{B}_m,\cdots,\bm{B}_1$ must collectively correspond to a configuration of directed edges within a $d\times (m+1)$ grid such that each directed edge connects only adjacent levels (i.e., neighboring columns), and every source node at the initial level can route information to all destination nodes at the $(m+1)$-th level.

The information transmission perspective offers valuable insights into the sparsity patterns of the factorization matrices used in OFT. As shown in Figure~\ref{fig:blk_diag}, $\bm{R}$ in the original OFT exhibits a block-diagonal arrangement. Although this configuration successfully reduces the parameter count, it cannot yield a densely populated matrix $\bm{R}$. Our objective is to construct a dense orthogonal matrix by combining $m$ sparse orthogonal matrices, striving to minimize the number of required trainable parameters. The information transmission view suggests two main criteria for this aim: \emph{(i)} \textbf{dense connectivity}, meaning that every node from the first layer maintains at least one route to each node in the final layer; and \emph{(ii)} \textbf{minimal free edges}, which requires minimizing the total edge count while respecting orthogonality. Notably, orthogonality introduces a precise limitation on connections across adjacent layers. Specifically, to ensure that each matrix $\bm{B}_i$ retains full rank (an essential property for orthogonality), there must be $d$ edges establishing a bijection between nodes at the $i$-th and $(i=1)$-th levels, meaning the minimum edge count between adjacent levels is $d$ (\emph{e.g.}, in Figure~\ref{fig:info_trans}, this is 4). Since these $d$ edges are mandated by orthogonality and their values are not learnable (\emph{e.g.}, in a $d\times d$ orthogonal matrix with $d$ nonzero elements, their values are necessarily $\pm 1$), they should not be included in the count of trainable edges. Owing to the parameter efficiency of orthogonal matrices compared to general dense matrices, orthogonality further imposes intricate dependencies among edges. For instance, in a $2\times 2$ orthogonal matrix, assigning zero to the $(1,1)$ entry necessitates the $(2,2)$ entry also being zero, illustrating how omitting a connection at one location can require the omission of another. Consequently, the set of valid edge patterns may be expanded or restricted by the orthogonality constraint. The information transmission approach proves particularly advantageous for analyzing OFT, as it focuses exclusively on the presence and arrangement of nonzero, trainable entries in $\bm{R}$, rather than their explicit values.

A straightforward fully connected configuration between two layers consists of $\mathcal{O}(d^2)$ connections (equivalent to using a full dense orthogonal matrix), which amounts to $d^2-d$ edges (for example, $d=4$ results in 12 connections). As demonstrated in Figure~\ref{fig:info_trans}, a possible matrix decomposition can be achieved with just $10$ connections, which is fewer than that required by a single dense orthogonal matrix. This perspective provides a graphical lens for analyzing the parameter efficiency of OFT, facilitating the identification of valid matrix decompositions. The design draws on the topology of butterfly graphs from the Cooley-Tukey algorithm~\cite{cooley1965algorithm}, known for enabling efficient dense connectivity between $d$ input and $d$ output nodes with only $\mathcal{O}(d\log d)$ connections. Specifically, the structure illustrated in Figure~\ref{fig:info_trans} employs $10$ connections for dense information flow, whereas a butterfly graph can achieve the same with merely $8$ edges. We proceed to explain how leveraging the butterfly topology enhances parameter efficiency.

\section{Orthogonal Parameterization by Butterfly Factorization}

The butterfly architecture has its roots in the Cooley-Tukey algorithm, where it is used to accelerate the fast Fourier transform. In the context of the Fourier transform, a modification in the frequency domain can induce widespread changes throughout the spatial domain, conceptually paralleling our scenario: every node at the initial layer should be able to relay information to all nodes at the terminal layer. Beyond its origins, the butterfly topology has also become a foundational structure in computer network design~\cite{solihin2015fundamentals,li2011linear} to enable scalable communication patterns. Assuming $k \geq 2$ and is a power of $2$, we define the \emph{butterfly factor} $\bm{B}^F(k)$ as follows:

\begin{equation}\label{eq:but_factor}
\begin{aligned}
    \bm{B}^F(k)=\begin{bmatrix}
    \text{diag}(\bm{d}_1) & \text{diag}(\bm{d}_2) \\
    \text{diag}(\bm{d}_3) & \text{diag}(\bm{d}_4)
    \end{bmatrix}\in\mathbb{R}^{k\times k},
\end{aligned}
\end{equation}

with each $\bm{d}_i\in\mathbb{R}^{\frac{k}{2}}$ for all $i$. For dimensions $d=2^N$, we recursively construct the $d$-dimensional \emph{butterfly matrix} $\bm{B}(d)\in\mathbb{R}^{d\times d}$ as:

\begin{equation}
\begin{aligned}
    \!\!\bm{B}(d)=\tilde{\bm{B}}(d,d)\!\cdot\!\begin{bmatrix}
    \bm{B}_1(\frac{d}{2})\!\!\!\!\! & \bm{0} \\
    \bm{0}\!\!\!\!\! &\bm{B}_2(\frac{d}{2})
    \end{bmatrix}= \tilde{\bm{B}}(d,d)\tilde{\bm{B}}(d,\frac{d}{2})\cdots\tilde{\bm{B}}(d,2),
\end{aligned}
\end{equation}

where $\bm{B}_1(\frac{d}{2})$ and $\bm{B}_2(\frac{d}{2})$ denote two butterfly matrices of dimension $\frac{d}{2}$. We introduce the \emph{butterfly component} as $\thickmuskip=2mu \medmuskip=2mu \tilde{\bm{B}}(d,k)=\text{diag}(\bm{B}^F_1(k),\cdots,\bm{B}^F_{d/k}(k))$, which forms a block-diagonal matrix of dimension $d\times d$ with blocks of size $k$, where each block $\bm{B}^F_i(k),\forall i$ corresponds to butterfly factors as specified in Equation~\ref{eq:but_factor}. With this construction, we can now parameterize an orthogonal matrix using butterfly matrices. This is accomplished by ensuring each factor $\tilde{\bm{B}}(d,k),\forall k$ that composes the butterfly matrix $\bm{B}(d)$ is itself orthogonal. Specifically, consider the block-diagonal matrix $\tilde{\bm{B}}(d,2)$ with blocks of size $2$; orthogonality for $\tilde{\bm{B}}(d,2)$ can be enforced by representing each $2\times 2$ block via a Cayley transform or $2$-dimensional rotation, consistent with \cite{liu2021orthogonal,qiu2023controlling}. The sparsity structure of each butterfly component is simply a permutation of that of $\tilde{\bm{B}}(d,2)$, facilitating orthogonal parameterization for all butterfly components. This construction yields an efficient orthogonal matrix parameterization composed of multiple $2\times 2$ orthogonal submatrices. Building on the work of \cite{chen2022pixelated}, we further extend this framework by defining the block butterfly component $\tilde{\bm{B}}^b(d,k)$, where every element of $\bm{d}_i$ is generalized to a $b\times b$ matrix. To preserve orthogonality in the block butterfly component $\tilde{\bm{B}}^b(d,2)$, each $2b\times 2b$ block is parameterized to be orthogonal. The structural sparsity pattern for other butterfly components $\tilde{\bm{B}}^b(d,k),k>2$ results from blockwise permutation of the pattern in $\tilde{\bm{B}}^b(d,2)$, which allows for orthogonal parameterization in a similar way. By aggregating these constructions, the BOFT forward computation is given by

\begin{equation}
    \begin{aligned}\nonumber
    \bm{z}=\big(\bm{R}(m,b)\cdot\bm{W}^0\big)^\top\bm{x},~~~~\text{s.t.}~\bigg{\{}\bm{R}(m,b)=\prod_{i=1}^m \tilde{\bm{B}}^b_{(d,i)}~~\&~~\underbrace{\big(\tilde{\bm{B}}^b_{(d,j)}\big)^\top\tilde{\bm{B}}^b_{(d,j)}=\tilde{\bm{B}}^b_{(d,j)}\big(\tilde{\bm{B}}^b_{(d,j)}\big)^\top\!=\bm{I}_d}_{\forall j\in [1, m]}\bigg{\}},
    \end{aligned}
\end{equation}

To streamline the notation, we define $\tilde{\bm{B}}^b(d,2^{m - i+1})$ as $\tilde{\bm{B}}^b_{(d,i)}$, where $\bm{I}_d$ denotes the $d$-dimensional identity matrix. The orthogonal transformation $\bm{R}(m,b)\in\mathbb{R}^{d\times d}$ arises as the sequential product of $m$ orthogonal butterfly matrices. For clarity, we denote BOFT parameterized with $\bm{R}(m,\frac{b}{2})$ as BOFT($m$,$b$), with $b \geq 2$. When $m=1$, BOFT($1$,$b$) coincides with the block-diagonal OFT described in~\cite{qiu2023controlling}, utilizing a block size of $b$. When BOFT($1$,$d$) is used, it is equivalent to the original OFT~\cite{qiu2023controlling} comprising an unconstrained, full orthogonal matrix. Furthermore, BOFT($\log \frac{2d}{b}$,$b$) employs the block butterfly matrix $\bm{B}^{\frac{b}{2}}(d)$ as $\bm{R}$, leading to a dense orthogonal matrix $\bm{R}$. More generally, BOFT($m$,$b$) involves $\frac{1}{2}(b-1)dm$ tunable parameters for adapting a linear layer of dimensions $d\times n$. If the butterfly matrix is adopted, that is $m=\log d,~b=2$, the parameter complexity of BOFT becomes $\mathcal{O}(d\log d)$. In contrast, a full dense orthogonal matrix as in the original OFT requires $\mathcal{O}(d^2)$ parameters, and a block-diagonal OFT partitioned into $r$ blocks uses $\mathcal{O}(bd)$. Thus, to yield a dense orthogonal matrix, the original OFT must set $b=d$, while BOFT can realize dense orthogonalization for any chosen $b$.

\textbf{Identity initialization in BOFT.} Standard finetuning protocols typically initialize the modified model to precisely replicate the pretrained weights at the outset, limiting drift from the original. For instance, LoRA adopts zero initialization for low-rank parameters. Similarly, BOFT initializes all butterfly factors as identity matrices (equivalently, initializing skew-symmetric matrices as zero for the Cayley transform).

\textbf{Multiplicative Dropout in BOFT.} The LoRA approach~\cite{hulora2022} incorporates a Dropout mechanism on its low-rank updates to avoid overfitting, and the classical Dropout~\cite{srivastava2014dropout} applies straightforwardly in that context. However, since BOFT updates weights multiplicatively, direct application of standard Dropout is unsuitable. As a remedy, we introduce a multiplicative Dropout tailored for BOFT: since $\bm{R}(m,b)$ is structured from $m$ orthogonal butterfly factors that can be reordered as $2b\times 2b$ block-diagonal orthogonal matrices, multiplicative Dropout operates by randomly selecting a proportion $p_1$ of butterfly matrices and, within each, a proportion $p_2$ of diagonal sub-blocks, which are then substituted with identity matrices.

\section{Intriguing Insights and Discussions}

\textbf{Expressivity of BOFT}. The butterfly-based structure, coupled with permutations, is capable of exactly reproducing many well-known efficient linear transforms~\cite{dao2019learning,dao2019kaleidoscope} (\emph{e.g.}, fast Fourier transform, Hadamard transform). However, the capability of our orthogonal butterfly matrix to closely approximate an arbitrary orthogonal matrix has not yet been thoroughly investigated. To address this, we perform a simulation experiment aimed at approximating a randomly generated dense orthogonal matrix~\cite{anderson1987generation} of dimension $512\times 512$. Figure~\ref{fig:para_eff} presents results averaged over 10 different seeds. The y-axis measures the approximation error, while the x-axis captures the count of effective trainable parameters. Each color-coded curve groups BOFT instances with a fixed block size, and for each, the point at the far left represents the error for BOFT($1$,$b$), corresponding to the original block-diagonal OFT setup with block size $b$. In general, BOFT consistently demonstrates superior parameter efficiency compared to OFT. For instance, BOFT($9$,$2$) achieves greater expressivity with significantly fewer parameters than BOFT($1$,$16$). Moreover, configurations with smaller $b$ and larger $m$ within BOFT typically realize better parameter efficiency: BOFT($6$,$4$), for example, matches the approximation error of BOFT($2$,$16$) using substantially fewer parameters. Essentially, the butterfly matrix structure forms a more constrained subset of the orthogonal group relative to the block-diagonal design, thus making BOFT provably more expressive than OFT given an equivalent block size.

\begin{theorem}[Expressivity of BOFT]\label{thm:exp_boft}
    BOFT possesses a strictly greater level of expressivity compared to OFT for any fixed block size. Specifically, to fully span all orthogonal matrices of size $d$, one can compose a sequence of butterfly matrices as $\bm{B}_{d-1,1}(d)\bm{B}^\top_{d-1,2}(d)\cdots\bm{B}_{1,1}(d)\bm{B}^\top_{1,2}(d)$, where every $\bm{B}_{i,j}(d)$ is a butterfly matrix.
\end{theorem}

This result, as formalized in Theorem~\ref{thm:exp_boft}, leads to a natural extension for BOFT in which the resultant orthogonal matrix takes the form $\thickmuskip=2mu \medmuskip=2mu \bm{R}^G(m_1,b_1,m_2,b_2,l)=\bm{R}_{l,1}(m_1,b_1)\bm{R}_{l,2}^T(m_2,b_2)\cdots\bm{R}_{1,1}(m_1,b_1)\bm{R}_{1,2}^T(m_2,b_2)$, where $\bm{R}_{i,j}^T(m,b)$ refers to an orthogonal matrix within BOFT. By choosing $m_1=m_2=\log d$, $b_1=b_2=2$, and $l=d-1$, $\bm{R}^G(m_1,b_1,m_2,b_2,l)$ becomes capable of spanning the entire orthogonal group, a property known as the kaleidoscope hierarchy~\cite{dao2019kaleidoscope}. Nevertheless, it is important to recognize that enhanced expressivity does not universally yield better performance in practical finetuning scenarios—indeed, fully finetuned models, though maximally expressive, often perform suboptimally. Effectively balancing expressivity and regularity is critical for strong model generalization. BOFT broadens the parameter space used in finetuning while imposing structural priors, thereby helping identify an optimal balance between these competing properties.

\textbf{Spectral properties}. Orthogonal finetuning typically achieves superior spectral characteristics compared to LoRA, as it ensures exact preservation of the spectral norm of the pretrained weight matrix $\bm{W}^0$. This outcome can be elucidated via singular value decomposition: $\bm{W}^0=\bm{U}\bm{\Sigma}\bm{V}^\top$, where $\bm{U}$ and $\bm{V}$ denote orthogonal matrices and $\bm{\Sigma}$ is a diagonal matrix containing the singular values. When OFT or BOFT are applied, they left-multiply an orthogonal matrix $\bm{R}$ to yield finetuned weights of the form $\bm{R}\bm{U}\bm{\Sigma}\bm{V}^\top$. This transformation leaves the largest singular value, i.e., the spectral norm of $\bm{W}^0$, unchanged. Maintaining the spectral norm in this way has been empirically linked to enhanced training stability and improved generalization~\cite{miyato2018spectral,yoshida2017spectral}. Further mathematical aspects are discussed in Appendix~\ref{app:sn_property}.

\textbf{Orthogonal finetuning as learning bilinear similarity}. The BOFT technique can be reinterpreted as the task of learning a bilinear similarity, specifically in the form $\bm{w}^0_i\bm{R}\bm{x}$, where $\bm{w}^0_i$ stands for the $i$-th neuron (or column vector) of the original weight matrix $\bm{W}^0$. Under this perspective, BOFT seeks to optimize the similarity matrix $\bm{R}$, while imposing a strict orthogonality constraint, closely relating it to approaches in distance metric learning~\cite{xing2002distance} and bilinear forms~\cite{roman2005advanced}.

\textbf{Inductive bias and generalization in BOFT}. In BOFT, $\bm{R}(m,b)$ generally describes a highly structured subgroup of the orthogonal group, restricting the hypothesis space and thereby introducing a specific inductive bias. We posit that the inductive bias imparted through butterfly factorization enhances generalization ability, owing to its structural similarities with various standard linear transforms, such as the discrete Fourier transform, discrete sine/cosine transforms, and the Hadamard transform~\cite{dao2019kaleidoscope}. Additionally, the sparse nature of the matrix factorizations used in BOFT can contribute further implicit inductive biases~\cite{gunasekar2017implicit,li2018algorithmic,liu2019neural}.

\textbf{Comparison to butterfly-based sparse training}. A number of previous studies~\cite{dao2019learning,dao2019kaleidoscope,chen2022pixelated,dao2022monarch} have explored sparse training techniques that leverage the butterfly parameterization. These works predominantly focus on direct reparameterization of weight matrices via butterfly parameterizations, training neural networks from their initial state. In contrast, \cite{dao2022monarch} investigates fine-tuning of pretrained models by projecting existing weights onto a variant of butterfly matrices, followed by optimization of the resulting projected components for downstream applications. The BOFT approach diverges significantly by introducing a fine-tuning methodology that applies shared transformation matrices across layers, rather than layer-wise butterfly reparameterization.

\input{rebuttal-results}

\section{Concluding Remarks and Limitations}

This work introduces Orthogonal Butterfly, a general-purpose parameter-efficient fine-tuning technique for foundation models rooted in the butterfly matrix framework. The main advancement in parameter efficiency is achieved by expressing a dense orthogonal matrix as a product of several sparse orthogonal matrices. To systematically discover viable matrix decompositions, we devise a graphical framework grounded in information transmission principles. Within this framework, we establish that butterfly structures are highly suitable for sparse orthogonal matrix factorization. We provide empirical evidence showcasing BOFT’s effectiveness for fine-tuning across diverse architectures, including large language models, vision models, and text-to-image generators. Our findings further demonstrate BOFT's competitive edge as a universally applicable model fine-tuning strategy.

Notwithstanding its promising performance, BOFT is not without drawbacks. The necessity to represent the final orthogonal transformation as a product of multiple orthogonal matrices introduces a moderate increase in training runtime compared to OFT. Optimizing BOFT’s computational efficiency during training remains an open challenge. Nonetheless, a notable benefit is that, post-finetuning, the orthogonal matrices learned by BOFT can be fused directly into the original model, thereby avoiding any additional inference time overhead. Furthermore, it is not yet clear whether the butterfly configuration represents the most optimal structure for information transmission. Our proposed framework, based on information transmission, invites the incorporation of insights from the field of computer networking, where analysis of network topologies for information dissemination is extensively studied. We anticipate that future research may uncover alternative, more efficient structures for constructing dense orthogonal matrices.